% Transformer text for reports/report_body.tex (PR #17), written for issue #10 by the transformer owner.
% Each block says which PENDING marker it replaces. Every number comes from a file in the repository:
%   dev:  results/transformer_comparison.csv, results/transformer_pairwise.csv, results/transformer_lr_search_ep6.csv
%   test: results/transformer_final_test/ (NOT YET RUN: the test rows below stay PENDING until the team approves the run;
%         protocol in docs/transformer_final_evaluation_protocol.md). Do not fill them in by hand.
% Citation keys follow report_body.tex (devlin2019, conneau2020).

% ---------------------------------------------------------------- Abstract / Introduction (one sentence each)
% Replaces: [PENDING: add one sentence on the transformer and LLM results if they are completed.]
Fine-tuned multilingual transformers were also compared: XLM-R reached 0.627 development weighted F1 (mean of three seeds)
against 0.610 for mBERT, while mBERT did not improve on the TF-IDF and BiLSTM baselines.
% After the test run, add: "On the test set, XLM-R and mBERT reached [X] and [Y] weighted F1."

% ---------------------------------------------------------------- Methodology
% Replaces the transformer part of: [PENDING: describe the fine-tuned mBERT, XLM-R and MuRIL/IndicBERT models ...]
\subsection{Fine-tuned transformers}
We fine-tune two pretrained multilingual encoders with a classification head on the three-class task: mBERT
(\texttt{bert-base-multilingual-cased}) \cite{devlin2019} and XLM-R (\texttt{xlm-roberta-base}) \cite{conneau2020}. Both use
one shared training script, so the pretrained model is the only difference besides the learning rate. Tweets use the
\emph{clean} text version; spellings are not normalised, because subword tokenizers already split unseen variants into known
pieces. The maximum input length is derived from the training set per tokenizer (the longest training tweet in subwords,
rounded up to a multiple of 8: 80 for mBERT, 96 for XLM-R), so no training or development tweet is truncated; batches are
padded dynamically. Training uses AdamW (weight decay 0.01, none on bias and LayerNorm weights), a linear warm-up over 10\% of
the steps followed by linear decay, batch size 32, gradient clipping at 1.0 and mixed precision, for at most 6 epochs with
early stopping (patience 2) on development weighted F1; the epoch with the best development weighted F1 is kept.

The learning rate is chosen per model on the development set from $\{1, 2, 3\}\times10^{-5}$ with one seed (42); the chosen
rate (mBERT $2\times10^{-5}$, XLM-R $1\times10^{-5}$) is then run with seeds 43 and 44, giving three seeds per model. An
initial search with 4 epochs was discarded because every run peaked in its last epoch; with 6 epochs most runs peak earlier.
On an RTX 4050 laptop GPU (6 GB) one epoch takes about 80~s for mBERT and 200~s for XLM-R.

% ---------------------------------------------------------------- Development results
% Add these rows to Table~\ref{tab:dev} (after the BiLSTM rows):
mBERT (3 seeds)  & 0.610 $\pm$ 0.005 & 0.614 $\pm$ 0.005 & 0.610 $\pm$ 0.005 \\
XLM-R (3 seeds)  & 0.627 $\pm$ 0.003 & 0.632 $\pm$ 0.003 & 0.631 $\pm$ 0.004 \\
% and extend the caption: "Transformer rows: results/transformer_comparison.csv; learning rate and best epoch chosen on dev."
% Replace the first sentence after "Table~\ref{tab:dev} shows development results." with:
% "All non-transformer models are within about 0.006 weighted F1 of each other."

% Paragraph to add after Table~\ref{tab:dev}:
XLM-R scores higher than mBERT on every seed (by 0.014 to 0.020 weighted F1), and it is about 0.015 above the best TF-IDF
model and the BiLSTM. The seed-to-seed spread is small (sample standard deviation 0.003 for XLM-R, 0.005 for mBERT), but the
gap is close to the sampling error of the development set: a paired bootstrap over the same development tweets (1{,}000
resamples, the median-seed run of each model) gives a difference of $-0.014$ with a 95\% interval from $-0.030$ to $+0.003$
for mBERT minus XLM-R. mBERT is level with the baselines. Both transformers find the neutral class hardest; XLM-R's gain
comes mainly from the positive class (recall 0.74 against 0.66 for mBERT). The learning rate was chosen with seed 42, so its
score is slightly optimistic; the mean over seeds 43 and 44 alone is 0.625 for XLM-R and 0.608 for mBERT.

% ---------------------------------------------------------------- Test results
% Replaces the transformer part of: [PENDING: add rows for the transformer and LLM systems.]
% NOT YET AVAILABLE. After `python -m src.evaluate_transformer_test --confirm-final-test`, copy the mean +/- sd values from
% results/transformer_final_test/transformer_final_test_report.md (section 1 and 2) into rows like:
% mBERT (3 seeds) & [wF1] $\pm$ [sd] & [macro F1] $\pm$ [sd] & [acc] $\pm$ [sd] \\
% XLM-R (3 seeds) & [wF1] $\pm$ [sd] & [macro F1] $\pm$ [sd] & [acc] $\pm$ [sd] \\
% (the BiLSTM table's first column header "BiLSTM text version" should then become "System").
% Sentence to keep with the rows: "The transformer checkpoints, learning rates and inference settings were frozen and
% hash-recorded before the test labels were read (docs/transformer_final_evaluation_protocol.md); all six runs are reported."

% ---------------------------------------------------------------- Conclusion and limitations
% Replaces the transformer part of: [PENDING: state what was and was not done for the transformers and the LLM.]
For the transformers, only mBERT and XLM-R base models were fine-tuned, on the \emph{clean} text version only. The learning
rate search used a single seed and three values, and other settings (batch size, warm-up, weight decay) were not tuned.
Learning rate and best epoch were selected on the development set, so development scores are optimistic. Larger models
(XLM-R large) and Indic-specific models were not part of this comparison.
% MuRIL/IndicBERT (issue #7) and the zero-shot LLM (issue #8) are covered by their owners.
